\begin{document}
\chapter{Spatial analysis and autocorrelation}
\section*{Overview}
This is the second session in the `autocorrelation' series, in which
we're going to cover the kinds of problems that can arise when we
collect data at different points in space. As with time series
analysis, it can be difficult to think of how we could collect data
\emph{outside} of space, so perhaps it's best to think of these
techniques as important for data where your explanatory variables
change across space. We will only be covering two-dimensional space in
this lecture series, and it is very rare for biologists (or really any
scientists) to need to control for spatial autocorrelation in a
three-dimensional setting. Typically, if you have variation in height
(elevation) in your data, this is fitted as an explanatory
variable. Perhaps the only exception to this in a biological setting
is marine data, and while I have only rarely seen studies of pelagic
data that need to account for all three dimensions, in such cases the
extension of these techniques (if not these precise packages) to such
data is trivial. To begin with, we're going to extend the classic
concept of correlation to account for \emph{spatial} autocorrelation,
which will give us a way to detect whether spatial autocorrelation
exists in our datasets. We'll then cover two common methods for
working with spatial data (autoregressive models and Generalized Least
Squares (GLS) models), and then finish up with some hands-on
experience working with spatial data.

\section{Detecting spatial autocorrelation}
While we have spent a great deal of time in this course talking about
the square of the correlation coefficient ($r^2$; see equation
\ref{r2-equ}), but we have not considered its square-root in much
detail. You have encountered the \emph{correlation coefficient} in
previous statistical classes, but let's review it here as we can
detect spatial autocorrelation using an extension of it. The
correlation coefficient, $r$, measures whether or not two vectors of
data are correlated. There are many different definitions of it but,
importantly, the \emph{Pearson's correlation coefficient} we are
focusing on today doesn't care which of them is the response or
explanatory variable---they're all the same to it. It is defined as:

\begin{equation}
  r = \frac{\text{Outliers in x} \cdot \text{Outliers in y}}{\text{Variance in x and y}} = \frac{\sum((x - \bar{x})(y-\bar{y}))}{\sqrt{\sum(x-\bar{x})^2}\sqrt{\sum(y-\bar{y})^2}}
  \label{pearson-equ}
\end{equation}

Where $x$ and $y$ are the two variables we are testing for
correlation, and $\bar{x}$ and $\bar{y}$ are the means of these two
variables. The idea of a correlation test is to see whether values
that are greater than the mean in one variable are paired with values
that are above or below the mean in the other variable (that's our
\emph{observation}, and our \emph{expectation} is they're not), in the
context of how much background \emph{variation} there is in the two
variables\footnote{Surprise! It's another test statistic!}. If there
were lots of unusually large values in $x$ that were paired with
unusually low values in $y$, we would tend to be multiplying lots of
large positive values with lots of negative values, and would sum
these up to get a very large negative value. This would be divided by
the overall variation in the two variables, and we would get a
negative number that was close to $-1$. The correlation coefficient
can only range between $-1$ and $+1$: values closer to $-1$ mean there
is a \emph{negative} correlation between the values, values closer to
$+1$ mean a positive value, and values closer to $0$ tend to mean
there's little correlation between the variables.

So how can we use this to measure the correlation of values \emph{in
  the same variable} through space? Using something called
\emph{Moran's I}, which is remarkably similar to a Pearson's
correlation coefficient, and is defined as:

\begin{equation}
  I = \text{How many things} \cdot \frac{\text{Outliers here} \cdot \text{Outliers there}}{\text{Variation}} = \frac{n}{\sum w_{i,j}} \cdot \frac{\sum w_{i,j}(x_i-\bar{x})(x_j-\bar{x})}{\sum(x_i-\bar{x})^2}
\end{equation}

Where $x$ is our only variable, $\bar{x}$ is still the mean of $x$,
$n$ is our number of data points, $w$ is a matrix (table) of
\emph{spatial weights} among the points and has dimensions $i$ (rows)
and $j$ (columns) each corresponding to a particular observation of
$x$ such that it has $n$ columns and $n$ rows. The spatial weights,
$w$, are directly proportional to the \emph{distances} among points:
closer points are weighted higher (you can decide how much higher; see
GLS below). Thus Moran's I is just the same, really, as a Pearson's
correlation coefficient, only now instead of looking to see whether
outliers in $x$ and $y$ are correlated, we're looking to see whether
outliers in $x$ are in roughly the same area in space. In order to
bring that concept of space into this, we need to weight our points
according to how close they are, as otherwise we would simply be
asking how correlated all points within a variable are with each
other, which isn't useful for our purposes\footnote{Indeed, it would
  just be the \emph{variance} of $x$}. There are two (equally fine)
way to view the standardization term $\frac{n}{\sum w_{i,j}}$ in all
of this: you can view it as dividing through by
$\frac{\sum w_{i,j}}{n}$, in which case it's part of the
\emph{variance} component (take the overall variation in the weights
matrix given how many points we have), or you can view it as you
viewed the degrees of freedom in $F$-ratios, in which case we're
(again) controlling for how many observations we have and also where
they are in space. Either is fine.

Moran's I values can be compared with an expected distribution of
Moran's I values if there were no aggregation of similar values across
space in exactly the same way as any other test statistics. It's quite
common to generate a null distribution of I values through some kind
of simulation procedure, however, because there's a dependence on the
locations of your data points in space that can be tricky to solve
analytically. Either way, the end result is the same: you get a
distribution of expected Moran's I values, and you compare your
observed value with that to find out what fraction of the distribution
is greater than your observation\footnote{If you have ever taken a
  course in bootstrapping, you might be a bit concerned that I've
  glossed over quite a lot of details about how to compare an observed
  value to a bootstrapped distribution. Yet you might also note that,
  from the definitions of frequentist distributions I gave you earlier
  in the class, frequentists treat probability as a set of
  observations in the long-run and so there's no difference really
  between bootstrapping and analytical expectations. If you haven't
  taken a class in bootstrapping, move on! Nothing to see
  here...}. Note that you can also have spatial disaggregation: values
of Moran's I that are less than $0$ and reflect points that are close
to one-another in space being different from each other.

\section{Accounting for spatial autocorrelation}
Imagine you've calculated Moran's I and it appears to be significantly
different from zero: you've got spatial autocorrelation of some
kind. Or, perhaps, you have a strong \emph{a priori} reason to suppose
you've got spatial autocorrelation and you don't want to test for it
because you know it's there\footnote{Good for you! You have been
  paying attention!}. What now? There is a huge body of literature
devoted to modeling spatial data, and I'm going to present to you two
reasonably straightforward methods for dealing with it. If you are
interested, I cannot recommend Dormann \emph{et al.} (2007) in
\emph{Ecography} highly enough. It's showing its age (it tests its
models on a Pentium 4!), but the general concepts involved really
haven't changed that much with the exception of recent focuses on
Wavelet and spatial point-pattern analysis\footnote{Which really,
  truly are beyond the scope of this course, but I can point you in
  the direction of good introductions if you are interested.}.

\subsection{Simultaneous AutoRegressive models}
\emph{Simultaneous AutoRegressive} models (SARs) are generalizations
of the autoregressive models we met last time in a temporal
context. The ``simultaneous'' component here accounts for the fact
that we can regress the data against nearby values of the response
variable, the explanatory variable, or both during model-fitting. You
pick between those options depending on the kind of autocorrelation
problem you think, or test to find, that you have. If you think the
underlying variable that is driving the autocorrelation is in the
explanatory variable (\emph{e.g.}, temperature is driving mood but
temperature is spatially autocorrelated) then you would regress your
data against nearby temperature values. If you think it's in the
response variable, (\emph{e.g.}, happy people make others around them
happy), then you might regress against nearby response variables. If
you're not sure... Go for both.

\subsection{Generalized Least Squares}
We met \emph{Generalized Least Squares} (GLS) models last time as a way to
flexibly add any kind of ARIMA structure into a model. The good news
is they can work in exactly the same way for spatial data. Where this
can be useful is if you found the decision about what kind of SAR
error structure to fit uncomfortable in the last section: with GLS,
you have a bit more flexibility, and can also have autocorrelation in
your \emph{error} component. This is analogous to the decision we made
in the time series session about whether we wanted to account for
autocorrelation of the response variable (autoregressive models) or
the error underlying the response variable (moving average models),
although the specifics are a little different.

To work with GLS models, you need to specify the functional form of
the autocorrelation you think is driving your data. For example, do
you think there's an exponential (decaying) effect of space, where
nearby points matter a lot and far away points matter very little?
Then you can specify an \emph{exponential} model of the form
$e^{-\frac{a}{d}}$, where $e$ is the exponential function, $d$ is the
\emph{distance} among points, and $a$ is an estimated parameter that
determines how fast the influence of nearby points decays. This is
exactly analogous to the moving average models from last time, only
now instead of an order (first, second, third, etc.) determining how
many timesteps we ``listen'' to from earlier in the time series, now
we have a continuous $d$istance and an amount we listen (the
exponential function with its slope determined by $a$).

\section{Hands on with spatial data in \R}
I'm not going to sugar-coat it: the hardest part of spatial analysis
in \R is the actual \R. For lots of very good computational
reasons\footnote{Most of which involve the fact that you have to work
  with spatial distance matrices, because matrices tend to take up
  lots of memory and are slow to manipulate} spatial analysis in \R
can be a real pain. So please view what follows as a cookbook: when
you come to having to do your own spatial analysis in \R, take these
instructions and modify them to suit your own purposes.

First off, as always, I'm going to simulate some data. This requires
the use of a rather fancy package to generate spatially-autocorrelated
errors. I'm happy to give a very cursory introduction to how the
package works if you're interested, but please just treat it as a
magic black box unless you're willing to (1) take a serious class in
spatial analysis\footnote{...and I do mean a \emph{serious} class in
  spatial analysis} or (2) read the vignette of the package
carefully. I'm constrained in this class by how much I can teach in
half a semester: using the analysis tools I'm teaching you in a real
analysis is fine, but using these simulation tools without
understanding them more deeply could lead to Very Bad Things.

\begin{minted}{splus}
  library(RandomFields)
  locations <- expand.grid(seq(-1,1,.1), seq(-1,1,.1))
  error <- matrix(RFsimulate(
      RMexp(), x=locations[,1], y=locations[,2])@data[,1],
    nrow=length(unique(locations[,1])),
    byrow=TRUE)
  temperature <- matrix(RFsimulate(
      RMexp(), x=locations[,1], y=locations[,2])@data[,1],
    nrow=length(unique(locations[,1])),
    byrow=TRUE)
  data <- data.frame(
    error=as.numeric(error), temperature=as.numeric(temperature),
    y=as.numeric(row(error)), x=as.numeric(col(error))
    )
  data$mood <- with(data, 3 + temperature*.5 + error)
\end{minted}%$

We now have a \texttt{data.frame}, called \texttt{data}, that contains
some information about people's average \texttt{mood}s through space
as a function of \texttt{temperature}. In our simulated data people
are strange, and seem to prefer a warmer environment. OK, so how do we
go about telling \R that we have spatial data, and then how do we use
Moran's I to test whether there's spatial autocorrelation in this
dataset?

\begin{minted}{splus}
  # Load a useful package
  library(spdep)
  # Copy our data
  sp.data <- data
  # Tell R our new data structure is spatial
  coordinates(sp.data) <- ~x+y
  # Build a spatial weight matrix
  neighbors <- knn2nb(knearneigh(sp.data, k=8))
  weights <- nb2listw(neighbors)
  moran.test(data$mood, weights)
  # For good measure, check our residuals for autocorrelation
  naive.model <- lm(mood ~ temperature, data=data)
  moran.test(residuals(naive.model), weights)
  # ...and this has the consequence that our coefficients are wrong
  summary(naive.model)
  # ... the slope should be 0.5 and the intercept 3
\end{minted}%$

A few things about what \R is doing ``behind the scenes'' here. By
giving (one) of our datasets \texttt{coordinates}, we're implicitly
telling \R to create a spatial data object for us. This is, for
reasons I alluded to above, a quite tricky thing to do, and if you
have a very large dataset then make sure you have a very large
computer. Second, we build our weights matrix using (in this case) the
eight nearest neighbors of each point. \R can do a lot with that, and
it might surprise you that this is enough to detect autocorrelation:
you can use more neighbors if you want, but remember that the
importance of space tails off with distance and setting this number
too high can take up a \emph{lot} of memory very quickly.

Most importantly, from the perspective of the statistics, we can see
that both our response variable and our residuals in our model show
spatial autocorrelation. This is a problem: it means our data, and our
inferences from that data, are pseudoreplicated, and so we need to
account for this in our model. So let's fit an SAR model to account
for this:

\begin{minted}{splus}
  # Fit the model
  sar.model <- errorsarlm(mood ~ temperature, data=data, weights)
  # Check the output (it's correct!)
  summary(sar.model)
  # Check for residual autocorrelation (there's none!)
  moran.test(residuals(sar.model), weights)
\end{minted}

Voil\'{a}! We've managed to get the correct parameter estimates out of
our model, and removed the residual autocorrelation from our model,
all with a SAR model. This particular kind of model only accounts for
autocorrelation in our errors, but this was still sufficient to get
accurate parameter estimates despite me simulating my explanatory
variable and my errors to have spatial autocorrelation. So the
take-home from this is not to worry too much if you don't know what
kind of SAR model to fit: one with autocorrelation in the errors is
normally sufficient. Let's move onto a GLS model to see how we fare
there...

\begin{minted}{splus}
  # Load the package
  library(nlme)
  # Fit the model (note we're using sp.data now)
  gls <- gls(mood ~ temperature, data=sp.data, corr=corExp())
  # Check the output - it's correct!
  # - p-values are rounded so "0" means "very small"
  summary(gls)
  # Check for residual autocorrelation (there's none!)
  moran.test(residuals(gls), weights)
\end{minted}

The nice thing about the GLS is we have a lot of flexibility about the
form of the autocorrelation we think is going on---there are many
options beyond \texttt{corExp} for exponential decay. The downside of
this flexibility is it can be mind-numbingly slow to fit a model. Even
in this reasonably small test-case, it took a very long time. There
are other packages that offer spatial GLS regression, and many of
those are much faster, but there is an advantage to having
\texttt{nlme::gls} in your back-pocket: it can fit mixed effects
models. Such models can be fit in Bayesian land, and there are plenty
of tutorials for \texttt{stan} and \texttt{BUGS} out
there\footnote{JAGS is not so spatial-friendly, sorry!}, but perhaps
we can talk about those some other time.

\section{Coda: space isn't always time, and a statistical fix isn't always the right answer}
I can't leave you without mentioning something that crops up a lot in
my field: the modeling of space as if it were evolutionary time. A
long time ago, in a galaxy far, far away, someone thought it would be
a very good idea to account for correlation through time on a
phylogenetic tree as if it were correlation through space [see
Diniz‐Filho \emph{et al.} (1998) in Evolution]. It was a very good
idea, and under some circumstances it can be useful\footnote{See, for
  example, a paper I thought I would hate being involved with but that
  actually changed my mind on a lot of these eigenvector approaches;
  Morales‐Castilla \emph{et al.} (2017) in Global Ecology and
  Biogeography.}, but in general it is absolutely not a good
idea. This was definitively shown quite some time ago [see Freckleton
\emph{et al.} (2011) in The American Naturalist], but every now and
again someone chances upon the idea of using some sort of spatial
method in a new context. If you are going to use spatial methods for
something that isn't space, do make sure they're mathematically simple
to describe\footnote{A single equation, with a handful of parameters,
  for the correlation structure.}, or else you will end up
over-engineering something.

Equally, and I might add that this applies to everything I have taught
you in this course, do remember that \emph{a statistical fix is no
  substitute for biological insight}. There are very few, if any,
cases where having all of the important variables that drive a system
in a model will go wrong in a straight-forward multiple regression of
the type you learned in lecture 4. Everything else, from hierarchical
models to spatial autocorrelation, is only needed if you don't have
all the right information to hand. A good example of this is the
Mantel test in spatial statistics, which I have deliberately not
taught you because it can so easily lead to problems in spatial
analyses [see Legendre \emph{et al.} (2015) in Methods in Ecology \&
Evolution---the answer to the title of their paper is
``no'']. Anything that ``cancels out'' or ``controls for'' something
inevitably means that there is some process that is being
ignored. Modeling spatial autocorrelation with a specific
autoregressive or GLS function is not ``canceling out'' the
autocorrelation, it's measuring its impact. An even better thing to
do, of course, would be to find out what is \emph{driving} that
autocorrelation and measure it.

\section{Exercises}
As ever, below is some code to load today's dataset into \R, and do a
few modifications to it. This dataset is a little different, in that
it's loaded into your \R session by the code you're going to execute
below. This code is very useful, not just because the dataset is
useful, but also because it shows you how to extract information about
particular points on the Earth's surface (in this case random ones). I
can assure you that there will come a day when you will be very
grateful for remembering this exercise...

\begin{minted}{splus}
  library(raster)
  r <- getData("worldclim",var="bio",res=10)
  points <- data.frame(expand.grid(seq(0,50), seq(-50,-150)))
  names(points) <- c("lat", "long")
  sp.points <- points
  coordinates(sp.points) <- ~long+lat
  data <- data.frame(extract(r, sp.points))
  names(data) <- c("temp.mean","diurnal.range", "isothermality",
    "temp.season","max.temp","min.temp","temp.range","temp.wettest",
    "temp.driest","temp.mean.warmest","temp.mean.coldest","precip",
    "precip.wettest","precip.driest","precip.season","precip.wettest",
    "precip.driest","precip.warmest","precip.coldest")
  data$lat <- points$lat
  data$long <- points$long
  data <- na.omit(data)
\end{minted}

\begin{enumerate}
\item Today's dataset is of global weather/climate conditions in
  roughly the present day. The code above shows you how to download
  that, and then extract information about these environmental
  variables at a set of points on the Earth's surface. You can adapt
  this code to study some particular points you're interested in if
  you wish!
  \begin{enumerate}
  \item Regress two of the environmental variables against one-another
    in a standard regression. What does it show?
  \item Is there significant spatial autocorrelation in either
    variable or the model residuals?
  \item Use either GLS regression or SAR regression to account for
    potential biases in the data. What does this show?
  \end{enumerate}

\item Today's code challenge is all about making it easier for you to
  grab climate information for any point on the planet's surface. Do
  pay attention to the code I have given you above...
\end{enumerate}
\begin{minted}{splus}
  grab.cliamte <- function(lat, long){
    raster <- getData(____, var="bio", res=10)
    points <- data.frame(lat=___, long=___)
    coordinates(____) <- ___ ___ ___ ___
    return(____(____, points))
  }
\end{minted}
